\subsection{Лекция 3 (3.10.26) МКO для решения уравнения Пуассона}
В тесте \cvar{poisson1-fvm} из файла \ename{poisson_fvm_test.cpp}
реализовано решение одномерного уравнения Пуассона с граничными условиями первого рода
(см. теорию в \secref{sec:FVM}).
Проводится расчёт на сгущающихся сетках с количеством ячеек от 10 до 1000
и расчитываются среднеквадратичные нормы отклонения полученного численного решения от точного.
Решения сохраняются в vtk-файлы \ename{poisson1_fvm_n={}.vtk}.
Отталкиваясь от этой реализации необходимо:
\begin{enumerate}

\item Написать аналогичный тест для двумерного уравнения в области $[1\times1]$.

\item Провести серию двумерных расчётов на сгущающихся сетках. Использовать структурированные и скошенные сетки.
      Визуализировать в Paraview полученное на этих сетках решение.
      Построить графики сеточной сходимости решения и определить порядок аппроксимации метода.

\item В тестовой программе производится сборка в цикле по ячейкам (\secref{sec:poisson_fvm_cellbased}).
      Следует переписать процедуру сборки в циклах по граням (\secref{sec:poisson_fvm_facebased}) и убедиться
      в идентичности результатов.

\item Вместо условий первого рода использовать условия второго рода для всех границ двумерной области.
      Убедится, что полученный ответ зависит от начального приближения с точностью до константы.
      Получить правильное решение явно задав значение в одной из точек коллокации.
\end{enumerate}

\paragraph{Рекомендации к программированию}
\begin{itemize}
\item Для выполнения п.1 необходимо по аналогии с классом \cvar{Problem1D} создать класс \cvar{Problem2D}.
      При этом в качестве аргумента конструктора у этого класса нужно использовать класс абстрактной
      двумерной сетки \cvar{IGrid2D}.
      Предложенный код уже не зависит от размерности задачи, поэтому в новом классе нужно определить только
      методы определения точного решения \cvar{exact_solution} и правой части \cvar{f}.
\item Необходимо создавать два вида сеток: структурированную и скошенную произвольную.
      Необходимые инструменты для работы с двумерными сетками описаны в \secref{sec:prog_grid}.
\item График сеточной сходимости следует строить по аналогии с графиком на \figref{fig:poisson_convergence}.
      в логарифмических осях, где по оси абсцисс отлежено разбиение, а по оси ординат -- норма.
      Разбиение $N$ -- это характерный линейный размер области расчёта, делённый на характерный линейный размер ячейки.
      Для получения корректного порядка аппроксимации в двух- и трёхмерных задачах следует внимательно отнестись к вычислению последнего параметра.
      Для двумерной/трёхмерной области характерный размер можно определить как квадратный/кубический корень от объёма.
\item Для программирования алгоритмов \cref{eq:fvm_assem_internal,eq:fvm_assem_bc1}, необходимых для выполнения п.3
      нужно получить список внутренних и граничных граней сетки. Это можно сделать используя методы
      \cvar{IGrid::internal_faces}, \cvar{IGrid::boundary_faces}. Поскольку в заданиях рассматриваются задачи
      с однотипными граничными условиями (только первого типа или только второго типа), то отдельно делить
      граничные грани по типам условий не нужно.
\item В п.4 для вычисления $q_s$ нужно путём аналитического интегрирования вычислить значения нормальной производной точного решения
      в центрах всех граничных граней. Следует внимательно обрабатывать направление $\vec n$, которое должно смотреть наружу области.
\begin{cppcode}
double q_face(size_t iface) const{
    Point p = grid_->face_center(iface);
    if (p.y < 1e-3){
        return ... // q(p) для нижней границы
    } else if (p.x > 1 - 1e-3){
        return ... // q(p) для левой границы
    } ...
}
\end{cppcode}
\item В чистой задаче Неймана решение определено с точностью до константы. В предложенном примере
      используется итерационный метод решения СЛАУ (Amgc). Поэтому решение будет зависеть от начального приближения.
      В исходной версии начальное приближение задаётся нулевым в строке
\begin{cppcode}
u_ = std::vector<double>(grid_->n_cells(), 0.0);
\end{cppcode}
      метода \cvar{BaseProblem::solve()}.
      Изменяя это приближение можно получать решения, смещённые на константу.
\item Чтобы получить единственное решение, не зависящее от начального приближения
      нужно явно определить значение $u$ в одной из ячеек. Например в первой:
      \begin{equation*}
          u_0 = u^e(\vec x_0).
      \end{equation*}
      Это выражение следует явно учесть в матрице левой части и столбце правой части после их сборок следующим образом:
\begin{cppcode}
mat.set_unit_row(0); // единица в диагональ, ноль вне диагонали
rhs[0] = exact_solution(grid_->cell_center(0));
\end{cppcode}
\end{itemize}
\label{sec:prog_grid}
